%!TEX root = ../template.tex
%%%%%%%%%%%%%%%%%%%%%%%%%%%%%%%%%%%%%%%%%%%%%%%%%%%%%%%%%%%%%%%%%%%
%% chapter1.tex
%% NOVA thesis document file
%%
%% Chapter with introduction
%%%%%%%%%%%%%%%%%%%%%%%%%%%%%%%%%%%%%%%%%%%%%%%%%%%%%%%%%%%%%%%%%%%

\typeout{NT FILE chapter1.tex}%

\chapter{Introduction}
\label{cha:introduction}

\glsaddall[types={main,symbols}]

\prependtographicspath{{\DIRFIGURES/Covers/}}

% epigraph configuration
\epigraphfontsize{\small\itshape}
\setlength\epigraphwidth{12.5cm}
\setlength\epigraphrule{0pt}


\includegraphics[height=2cm]{\DIRIMAGES/novathesis-insignia-outline}\hfill

\includegraphics[height=2cm]{\DIRIMAGES/novathesis-text}

\noindent This is the \gls{novathesis} \hologo{LaTeX} template \ntindex[template!version]{Version} \novathesisversion\ from \novathesisdate.

\epigraph{
  This work is licensed under the \href{https://www.latex-project.org/lppl/lppl-1-3c/}{\hologo{LaTeX} Project Public License v1.3c}.
  To view a copy of this \ntindex[template!license]{license}, visit the \href{https://www.latex-project.org/lppl/}{LaTeX project public license}.
}

% \section{Welcome to the \novathesis\ Template}
% \label{sec:if_you_use_this_template}

This first Chapter introduces the \gls{novathesis} template and how it is organized. In Chapter~\ref{cha:users_manual} you can find some specific instructions on how to use this template.  Chapter~\ref{cha:a_short_latex_tutorial_with_examples} shows some examples and give some hints on how to write your text. Please read these next Chapters carefully.

\section{The \emph{\texorpdfstring{\textbf{nova}thesis}{novathesis}} Template}
\label{sec:a_bit_of_history}

\ntindex[template!history]{}

The \gls{novathesis} template was born at the \gls{DI} of  \gls{FCT} of \gls{NOVA}, Portugal.  But the user base grew… initially grew to other Departments of NOVA FCT, then to other Schools of NOVA, and later to other Schools of other Universities.  Currently more than~25 Schools are natively supported by the \gls{novathesis} template.  The full list of supported Schools can be found at the
\href{https://github.com/joaomlourenco/novathesis}{Project's \ntindex[github]{GitHub} page}.

    
\subsection{Your Time is Precious}
\label{sub:time_is_money}

Did you learn how to drive by sitting by the wheel and throwing your car into the road?  Most probably you did take your time \emph{learning the rules} and \emph{practicing} first! Likewise, it is not wise to throw yourself at the task of writing a thesis/dissertation in \hologo{LaTeX} without seriously considering the following \ntindex{recommendation}!

\begin{highlightbox}
  If you are going to spend zillions of hours writing your thesis/dissertation using the \gls{novathesis} \hologo{LaTeX} template (or some other \hologo{LaTeX} template), be wise and spend a couple of hours learning how to use it properly by reading this manual.  And then, be even wiser, and spend a few more hours \href{https://github.com/joaomlourenco/novathesis/wiki\#learning-latex}{learning some \LaTeX}.  I am sure that the time you are investing now will pay itself countless times before you submit your thesis/dissertation.\\\parbox{\linewidth}{\raggedleft---~\emph{João Lourenço}}
\end{highlightbox}

\subsection{Recognition}
\label{sub:recognition}

\ntindex[recognition]{}

The \gls{novathesis} template was born in~1996, and what you see now accumulates to many many hundreds (thousands?!) of working hours, unpaid and stolen from family and friends.  This work is available to the community under the \href{LaTeX project public license}{\hologo{LaTeX} Project Public License v1.3c}, which means you are entitled to use it for free and change it at your will.  However, if you decide to use this template to write your thesis/dissertation, \textbf{be fair to the developers} and:
\begin{enumerate}
  \item \ntindex[novathesis!citation]{} Cite the \gls{novathesis} manual~\cite{novathesis-manual} in a place of your choice (e.g., in the \emph{Acknowledgments}) of your thesis/dissertation with “\verb!\cite{novathesis-manual}!” .  If you cite it this way, the correct entry will be added automatically to your bibliography (no need to worry with the necessary BibTeX entry, as it will be added automatically);
  \item Go to the
\href{https://github.com/joaomlourenco/novathesis}{\ntindex[github!project web page]{project web page} in \ntindex[github]{GitHub}} and give the project a \ntindex[github!stars]{star} (marked with a red ellipse at the top-right in Figure~\ref{fig:github}); and
  \item Make a \ntindex[donations]{donation} by visiting the \gls{novathesis} project page and clicking in the button marked with a green ellipse at the top-center in Figure~\ref{fig:github}).  Alternatively, just click \href{https://www.paypal.com/donate/?hosted_button_id=8WA8FRVMB78W8}{\fcolorbox{DarkGreen}{gray!15}{\textbf{~HERE~}}} and your browser will be directed to the right page.
\end{enumerate}

\begin{figure}[htbp]
    \centering
    \includegraphics[width=0.5\linewidth]{github1}
    \caption{The \gls{novathesis} project web page in GitHub.}
    \label{fig:github}
\end{figure}


\section{Getting Started}
\label{sec:getting_started}

The \gls{novathesis} template provides an \emph{easy to use} environment for writing your thesis. Detailed instructions on how to set up the template using \Overleaf\ or a local \hologo{LaTeX} installation are provided in the \emph{User's Manual} (see \autoref{sec:getting_started_manual}). This is particularly useful for students in \gls{STEM} fields who need to produce high-quality \gls{PDF} documents.


\section{Your Thesis Roadmap}
\label{sec:roadmap}

Writing a thesis is a journey. Here is the recommended workflow for using this template:

\begin{enumerate}
  \item \textbf{Configuration:} Start by editing the files in \ntindex[configuration!0-config/]{\latexinline{0-Config/}}. Select your school, set your document type (\ntindex[document types!phd]{PhD}, \ntindex[document types!msc]{MSc}, etc.), and fill in your metadata (title, author, supervisors) in \ntindex[configuration!0-config/3\_cover.tex]{\texttt{3\_cover.tex}}.
  \item \textbf{\ntindex[front matter]{Front Matter}:} Write your acknowledgments, dedication, and abstracts in the \ntindex[front matter!1-frontmatter/]{\texttt{1-FrontMatter/}} folder.
  \item \textbf{\ntindex[main matter]{Main Content}:} Create separate \texttt{.tex} files for each chapter in \ntindex[main matter!2-mainmatter/]{\texttt{2-MainMatter/}}. Register them in \ntindex[configuration!0-config/4\_files.tex]{\texttt{0-Config/4\_files.tex}} to include them in the build.
  \item \textbf{\ntindex[bibliography]{Bibliography}:} Keep your references in \texttt{4-Bibliography/bibliography.bib}. The template uses \ntindex[bibliography!biblatex]{\latexinline{biblatex}} to handle formatting automatically.
  \item \textbf{Build \& Iterate:} Use \ntindex[building!make]{\latexinline{make}} or your editor's build command frequently to check the results.
\end{enumerate}


\section{Getting Help}
\label{sec:getting_help}

\ntindex[help]{}

No! You don't have to use this template to write your thesis.  You don't even have to use \LaTeX.  However, writing a thesis is serious stuff, and which tool you shall use to write it is not a decision to make lighthearted.

\hologo{LaTeX} is hard enough by itself.  This template aims at making your life easier, but not easy. If you choose to use this template to write your thesis, you are very welcome.  However, don't expect me to provide you help with \LaTeX.  Look for help with your friends (you have some friends, don't you?), or search the web, or try even to read some book(s) on \LaTeX. In the end you will certainly find the experience rewarding.

When you come to the point of “\emph{How do I do this with the \novathesis\ template?}”, remember…

\begin{enumerate}
  \item To check the \href{https://github.com/joaomlourenco/novathesis/wiki}{\gls{novathesis} wiki} and have some hope!  \emojiSmile
  \item \href{https://www.google.com}{Google} is your best friend.
  \item Search the \href{https://github.com/joaomlourenco/novathesis/discussions}{GitHub Discussions page} for a question related to yours.  \emph{If and only if} you don't find one, then post your own question in English please!
  \item Search the \href{https://www.facebook.com/groups/novathesis}{NOVAtheis Facebook group} for a question related to yours.  \emph{If and only if} you don't find one, then post your own question in either Portuguese or English, at your preference.
\end{enumerate}

When you post your own question, remember to \textbf{always} state the \gls{novathesis} version number you are using and referring to.

\begin{warningbox}
	\centering
Please do not attempt to contact me directly (email, Messenger, etc)…\\I WILL NOT REPLY!
\end{warningbox}


\subsection{Suggestions, Bugs and Feature Requests} % (fold)
\label{sub:suggestions_bugs_and_feature_requests}

\begin{description}
  \item[Help:] If you just need some help, see above \Autoref{sec:getting_help}.
  \item[Suggestion:] \ntindex[suggestions]{Suggestions} Do you have a suggestion/recommendation? Please add it to the wiki and help other users!
  \item[Bug:] \ntindex[bugs]{Bugs} Did you find a bug? Please open an issue. Thanks!
  \item[New Feature:] \ntindex[feature requests]{Feature Requests} Would you like to request a new feature (or support of a new School)? Please open an issue. Thanks!

\end{description}



% subsection suggestions_bugs_and_feature_requests (end)




\section{Donors}
\label{sec:donations}

\ntindex[donations]{}

The \href{https://github.com/joaomlourenco/novathesis/wiki#donators}{list of \emph{Donors}} is available in the \gls{novathesis} Project page.


\section{Disclaimer}
\label{sec:disclaimer}

\ntindex[disclaimer]{}

Although the \gls{novathesis} template is endorsed by some Schools (e.g., \href{https://www.fct.unl.pt/estudante/informacao-academica/teses-e-dissertacoes}{linked from NOVA FCT web site}), the \gls{novathesis} template \textbf{this not an official template} for any School.

The \gls{novathesis} template exists to make your life easier and we do our best to make it compliant to the supported ($+25$) Schools' regulations but, in the end of the line, you and only you are accountable for both the look and the contents of the document you submit as your thesis/dissertation.
